 
 % !TEX program = pdflatex
% !TEX option  = -synctex=1 -interaction=nonstopmode -output-directory=build

% ==========================================================
% PLANTILLA BEAMER PRO
% Moderna, institucional y editable
% Compatible con pdfLaTeX
% ==========================================================

\newif\ifhandout
\handoutfalse          % <-- Presentación
%\handouttrue          % <-- Handout imprimible

\ifhandout
\documentclass[10pt,handout,aspectratio=169]{beamer}
\else
\documentclass[11pt,aspectratio=169]{beamer}
\fi

% ==========================================================
% PAQUETES
% ==========================================================
\usepackage[spanish,es-noshorthands]{babel}
\usepackage[T1]{fontenc}
\usepackage[utf8]{inputenc}
\usepackage{lmodern}
\usepackage{amsmath,amssymb,amsfonts}
\usepackage{mathtools}
\usepackage{graphicx}
\usepackage{xcolor}
\usepackage{tikz}
\usepackage{booktabs}
\usepackage{array}
\usepackage{multicol}
\usepackage{ragged2e}
\usepackage{hyperref}
\usepackage{enumitem}
\usepackage{pgfplots}
\usepackage{fontawesome5}

\pgfplotsset{compat=1.18}

% ==========================================================
% MODO HANDOUT
% ==========================================================
\ifhandout
\setbeameroption{hide notes}
\fi

% ==========================================================
% PALETA INSTITUCIONAL
% ==========================================================
% ---------- Opción UCM / Agronomía ----------
\definecolor{UCMGreen}{RGB}{18,112,84}
\definecolor{UCMGreenDark}{RGB}{10,74,56}
\definecolor{UCMGreenSoft}{RGB}{232,244,239}
\definecolor{UCMGold}{RGB}{180,145,76}
\definecolor{Ink}{RGB}{35,43,52}
\definecolor{SoftGray}{RGB}{245,247,250}

% ---------- Opción UST ----------
%\definecolor{UCMGreen}{RGB}{0,70,127}
%\definecolor{UCMGreenDark}{RGB}{0,45,85}
%\definecolor{UCMGreenSoft}{RGB}{232,240,248}
%\definecolor{UCMGold}{RGB}{0,153,117}
%\definecolor{Ink}{RGB}{35,43,52}
%\definecolor{SoftGray}{RGB}{245,247,250}

% ==========================================================
% TEMA BASE
% ==========================================================
\usetheme{Madrid}
\usecolortheme[named=UCMGreen]{structure}

\setbeamertemplate{navigation symbols}{}
\setbeamertemplate{blocks}[rounded][shadow=false]
\setbeamertemplate{itemize items}[circle]
\setbeamertemplate{enumerate items}[default]

% ==========================================================
% COLORES BEAMER
% ==========================================================
\setbeamercolor{normal text}{fg=Ink,bg=white}
\setbeamercolor{structure}{fg=UCMGreen}
\setbeamercolor{title}{fg=white,bg=UCMGreenDark}
\setbeamercolor{subtitle}{fg=white}
\setbeamercolor{frametitle}{fg=white,bg=UCMGreen}
\setbeamercolor{framesubtitle}{fg=white}
\setbeamercolor{author}{fg=white}
\setbeamercolor{date}{fg=white}
\setbeamercolor{institute}{fg=white}

\setbeamercolor{block title}{fg=white,bg=UCMGreen}
\setbeamercolor{block body}{fg=Ink,bg=UCMGreenSoft}

\setbeamercolor{block title alerted}{fg=white,bg=orange!80!black}
\setbeamercolor{block body alerted}{fg=Ink,bg=orange!8}

\setbeamercolor{block title example}{fg=white,bg=UCMGold!85!black}
\setbeamercolor{block body example}{fg=Ink,bg=yellow!8}

\setbeamercolor{section in toc}{fg=UCMGreen}
\setbeamercolor{subsection in toc}{fg=Ink}

% ==========================================================
% FUENTES Y ESPACIADO
% ==========================================================
\setbeamerfont{title}{series=\bfseries,size=\LARGE}
\setbeamerfont{subtitle}{size=\large}
\setbeamerfont{frametitle}{series=\bfseries,size=\large}
\setbeamerfont{block title}{series=\bfseries}
\setbeamerfont{footline}{size=\scriptsize}
\setbeamerfont{section title}{series=\bfseries,size=\Huge}

% ==========================================================
% HIPERVÍNCULOS
% ==========================================================
\hypersetup{
	colorlinks=true,
	linkcolor=UCMGreenDark,
	urlcolor=UCMGreenDark
}

% ==========================================================
% COMANDOS PEDAGÓGICOS
% ==========================================================
\newcommand{\pp}{\pause}

\newenvironment{stepitemize}
{\begin{itemize}[<+->,leftmargin=*,itemsep=0.55em]}
	{\end{itemize}}

\newenvironment{stepenumerate}
{\begin{enumerate}[<+->,leftmargin=*,itemsep=0.55em]}
	{\end{enumerate}}

\newcommand{\IdeaClave}[1]{
	\begin{alertblock}{\faLightbulb\ Idea clave}
		#1
	\end{alertblock}
}

\newcommand{\Ejemplo}[1]{
	\begin{exampleblock}{\faPen\ Ejemplo}
		#1
	\end{exampleblock}
}

\newcommand{\Actividad}[1]{
	\begin{block}{\faUsers\ Actividad}
		#1
	\end{block}
}

\newcommand{\Definicion}[1]{
	\begin{block}{\faBook\ Definición}
		#1
	\end{block}
}

\newcommand{\Objetivos}[1]{
	\begin{block}{\faBullseye\ Objetivos de aprendizaje}
		#1
	\end{block}
}

\newcommand{\Resumen}[1]{
	\begin{block}{\faClipboardCheck\ Síntesis}
		#1
	\end{block}
}

% ==========================================================
% PORTADA
% ==========================================================
\newcommand{\PortadaInstitucional}{
	\begin{frame}[plain]
		\begin{tikzpicture}[remember picture,overlay]
			% fondo
			\fill[SoftGray] (current page.south west) rectangle (current page.north east);
			
			% banda superior
			\fill[UCMGreenDark] (current page.north west) rectangle ([yshift=-2.1cm]current page.north east);
			
			% franja decorativa
			\fill[UCMGold] ([yshift=-2.1cm]current page.north west) rectangle ([yshift=-2.28cm]current page.north east);
			
			% panel lateral decorativo
			\fill[UCMGreen!12] (current page.south west) rectangle ([xshift=2.6cm]current page.north west);
			
			% círculos suaves
			\fill[UCMGreen!18] ([xshift=1.3cm,yshift=-1.5cm]current page.north west) circle (0.9cm);
			\fill[UCMGold!25] ([xshift=2.2cm,yshift=-3.8cm]current page.north west) circle (0.55cm);
			\fill[UCMGreen!10] ([xshift=-1.2cm,yshift=1.2cm]current page.south east) circle (1.4cm);
			
			% iconos
			\node[anchor=north west,text=white] at ([xshift=0.9cm,yshift=-0.7cm]current page.north west)
			{\fontsize{18}{18}\selectfont \faSeedling\ \faLeaf\ \faChartLine};
			
			% “logo” textual provisional
			\node[anchor=north east,text=white] at ([xshift=-0.8cm,yshift=-0.7cm]current page.north east)
			{\small \textbf{FACULTAD / ESCUELA / CARRERA}};
			
			% título
			\node[anchor=west,text=Ink] at ([xshift=3.2cm,yshift=-3.6cm]current page.north west) {
				\begin{minipage}{0.72\paperwidth}
					{\usebeamerfont{title}\color{UCMGreenDark}\inserttitle\par}
					\vspace{0.3cm}
					{\usebeamerfont{subtitle}\color{UCMGreen}\insertsubtitle\par}
				\end{minipage}
			};
			
			% datos
			\node[anchor=west,text=Ink] at ([xshift=3.2cm,yshift=-6.1cm]current page.north west) {
				\begin{minipage}{0.62\paperwidth}
					{\normalsize \textbf{\insertauthor}}\par\vspace{0.15cm}
					{\small \insertinstitute}\par\vspace{0.15cm}
					{\small \insertdate}
				\end{minipage}
			};
			
			% caja inferior
			\fill[white] ([xshift=3.2cm,yshift=1.1cm]current page.south west) rectangle ([xshift=-1.0cm,yshift=2.6cm]current page.south east);
			\node[anchor=west,text=Ink] at ([xshift=3.5cm,yshift=1.85cm]current page.south west)
			{\small \faGraduationCap\ Presentación académica moderna \quad \faChalkboardTeacher\ Clase universitaria \quad \faBookOpen\ Material editable};
		\end{tikzpicture}
	\end{frame}
}

% ==========================================================
% SEPARADORES DE SECCIÓN
% ==========================================================
\AtBeginSection[]{
	\begin{frame}[plain]
		\begin{tikzpicture}[remember picture,overlay]
			\fill[UCMGreenDark] (current page.south west) rectangle (current page.north east);
			\fill[UCMGold] (current page.south west) rectangle ([yshift=0.35cm]current page.north east);
			\node[text=white,align=center] at (current page.center){
				{\usebeamerfont{section title}\insertsection\par}
				\vspace{0.4cm}
				{\Large \faLeaf\ \ \faChartBar\ \ \faSeedling}
			};
		\end{tikzpicture}
	\end{frame}
}

% ==========================================================
% ENCABEZADO Y PIE
% ==========================================================
\setbeamertemplate{headline}{}

\setbeamertemplate{footline}{
	\leavevmode
	\hbox{
		\begin{beamercolorbox}[wd=.28\paperwidth,ht=2.9ex,dp=1.1ex,leftskip=1em]{author in head/foot}
			\color{white}\small \textbf{\insertshortauthor}
		\end{beamercolorbox}
		\begin{beamercolorbox}[wd=.50\paperwidth,ht=2.9ex,dp=1.1ex,center]{title in head/foot}
			\color{white}\small \insertshorttitle
		\end{beamercolorbox}
		\begin{beamercolorbox}[wd=.22\paperwidth,ht=2.9ex,dp=1.1ex,rightskip=1em plus1fil]{date in head/foot}
			\color{white}\small \insertframenumber/\inserttotalframenumber
		\end{beamercolorbox}
	}
}

\setbeamercolor{author in head/foot}{bg=UCMGreenDark,fg=white}
\setbeamercolor{title in head/foot}{bg=UCMGreen,fg=white}
\setbeamercolor{date in head/foot}{bg=UCMGreenDark,fg=white}

% ==========================================================
% TOC
% ==========================================================
\begin{document}
	
	% ==========================================================
	% DATOS EDITABLES
	% ==========================================================
	\title[Fundamentos de Matemática Básica]{Fundamentos de Matemática Básica}
	\subtitle{Potencias, notación científica y detección de patrones}
	\author[Luis González Alcaino]{Luis González Alcaino}
	\institute[UCM]{Universidad Católica del Maule \\ Curicó \\ Ingeniería de Ejecución Agrícola}
	\date{Año académico 2026}
	
	% ==========================================================
	% PORTADA
	% ==========================================================
	\PortadaInstitucional
	
	% ==========================================================
	% ÍNDICE
	% ==========================================================
	\begin{frame}{\faList\ Ruta de la clase}
		\tableofcontents
	\end{frame}
	
	% ==========================================================
	\section{Inicio}
	
	\begin{frame}{\faBullseye\ Propósito de la clase}
		\Objetivos{
			\begin{itemize}
				\item Aplicar el concepto de potencia y sus propiedades.
				\item Transformar números a notación científica.
				\item Resolver multiplicaciones y divisiones en notación científica.
				\item Reconocer patrones, relaciones e inconsistencias en cálculos y resultados.
			\end{itemize}
		}
	\end{frame}
	
	\begin{frame}{\faComments\ Activación de conocimientos previos}
		\Actividad{
			En parejas, discutan:
			\begin{itemize}
				\item ¿Dónde aparecen números muy grandes o muy pequeños en agronomía?
				\item ¿Qué ventajas tiene escribirlos en forma abreviada?
				\item ¿Qué errores son frecuentes al trabajar con potencias?
			\end{itemize}
		}
	\end{frame}
	
	% ==========================================================
	\section{Desarrollo conceptual}
	
	\begin{frame}{\faBook\ Potencia de un número real}
		\Definicion{
			Si $a$ es un número real y $n$ es un entero positivo, entonces:
			\[
			a^n = \underbrace{a\cdot a\cdot a\cdots a}_{n\text{ factores}}
			\]
			donde $a$ es la base y $n$ es el exponente.
		}
		\IdeaClave{
			La potencia representa una multiplicación repetida de una misma base.
		}
	\end{frame}
	
	\begin{frame}{\faPen\ Propiedades de las potencias}
		\begin{block}{Propiedades fundamentales}
			\begin{align*}
				a^m\cdot a^n &= a^{m+n} \\
				\dfrac{a^m}{a^n} &= a^{m-n}, \quad a\neq 0 \\
				(a^m)^n &= a^{mn} \\
				(ab)^n &= a^n b^n \\
				\left(\dfrac{a}{b}\right)^n &= \dfrac{a^n}{b^n}, \quad b\neq 0
			\end{align*}
		\end{block}
		
		\Ejemplo{
			\[
			2^3\cdot 2^5 = 2^{8}=256
			\]
		}
	\end{frame}
	
	\begin{frame}{\faFlask\ Contexto agronómico}
		\Ejemplo{
			Una muestra de suelo contiene aproximadamente
			\[
			3.2\times 10^5
			\]
			microorganismos por gramo.
			
			Si se analizan $4\times 10^2$ gramos, entonces el total estimado es:
			\[
			(3.2\times 10^5)(4\times 10^2)=12.8\times 10^7=1.28\times 10^8
			\]
			microorganismos.
		}
		\IdeaClave{
			En notación científica, se multiplican coeficientes por separado y se suman exponentes de base 10.
		}
	\end{frame}
	
	% ==========================================================
	\section{Trabajo guiado}
	
	\begin{frame}{\faChalkboardTeacher\ Ejemplo guiado}
		\Ejemplo{
			Simplificar:
			\[
			\dfrac{(2x^3y^2)^2}{4xy}
			\]
		}
		\pp
		
		\begin{block}{Paso 1}
			\[
			(2x^3y^2)^2 = 2^2x^6y^4 = 4x^6y^4
			\]
		\end{block}
		
		\pp
		\begin{block}{Paso 2}
			\[
			\dfrac{4x^6y^4}{4xy}=x^5y^3
			\]
		\end{block}
		
		\pp
		\IdeaClave{
			Primero se aplican las propiedades de potencias; luego se simplifican coeficientes y exponentes.
		}
	\end{frame}
	
	\begin{frame}{\faUsers\ Actividad de aula}
		\Actividad{
			Resuelve en grupos:
			\begin{enumerate}
				\item $\left(3\times 10^4\right)\left(2\times 10^3\right)$
				\item $\dfrac{6\times 10^8}{2\times 10^3}$
				\item $(5a^2b)^3$
				\item $\dfrac{x^7}{x^2}$
			\end{enumerate}
			
			Luego comparen resultados e identifiquen:
			\begin{itemize}
				\item patrones;
				\item posibles errores;
				\item formas de verificar el resultado.
			\end{itemize}
		}
	\end{frame}
	
	% ==========================================================
	\section{Cierre}
	
	\begin{frame}{\faClipboardCheck\ Síntesis final}
		\Resumen{
			\begin{itemize}
				\item Las potencias permiten representar multiplicaciones repetidas.
				\item Sus propiedades simplifican expresiones algebraicas.
				\item La notación científica facilita el trabajo con números extremos.
				\item Detectar patrones ayuda a validar procedimientos y resultados.
			\end{itemize}
		}
	\end{frame}
	
	\begin{frame}{\faQuestionCircle\ Ticket de salida}
		\begin{block}{Responde individualmente}
			\begin{enumerate}
				\item ¿Qué propiedad de potencias te resultó más útil hoy?
				\item Expresa $0.000045$ en notación científica.
				\item ¿Qué error evitarías al multiplicar potencias de base 10?
			\end{enumerate}
		\end{block}
	\end{frame}
	
	% ==========================================================
	% DIAPOSITIVA FINAL
	% ==========================================================
	\begin{frame}[plain]
		\begin{tikzpicture}[remember picture,overlay]
			\fill[UCMGreenDark] (current page.south west) rectangle (current page.north east);
			\node[text=white,align=center] at (current page.center){
				{\Huge\bfseries Gracias}\par
				\vspace{0.4cm}
				{\Large \faSeedling\ \faLeaf\ \faChartLine}\par
				\vspace{0.6cm}
				{\normalsize Luis González Alcaino}
			};
		\end{tikzpicture}
	\end{frame}
	
\end{document}